\section{Introduction}
%==============================================================================
Glass bottles, windows, and clear plastics defeat conventional segmentation
networks for a fundamental reason: a transparent surface has no stable
appearance of its own. Its colour, texture, and edges are functions of
\emph{what lies behind it}, not \emph{what it is} \cite{xie2020translab}. This
breaks the brightness-constancy assumption on which appearance-driven segmenters
implicitly rely, with tangible consequences: an autonomous agent cannot brake
for a glass door it has segmented as the corridor behind it, and a manipulator
cannot grasp a beaker it perceives as the wall.

A recurring and intuitively appealing remedy is to inject \emph{physics}.
Because refraction bends light according to Snell's law, the apparent motion of
the background seen through glass should not match rigid camera motion,
producing a measurable inter-frame inconsistency; and the two refractive
interfaces of glass should produce a characteristic double-edge frequency
response. If a network were handed these signals directly, the reasoning goes,
it would localise transparency from \emph{causes} rather than borrowed
appearance.

\textbf{This is an analysis paper, not a method paper.} Rather than propose the
physics cues as an advance, we treat them as a hypothesis to be tested and,
ultimately, refuted. We implement both cues---an Optical-Flow Consistency
Violation (\ofcv) gate and a Boundary Resonance Field (\brf) prior---inside a
single architecture, \spectra, built on a self-supervised backbone, and ask one
question under controlled conditions: \emph{do physics-motivated cues improve
transparent-object segmentation once a strong representation is present?} Our
answer is \emph{no, not measurably}, and we regard establishing this carefully
as the contribution.

\textbf{Headline finding.} In a five-variant $\times$ 30-epoch $\times$ 3-seed
ablation, removing or adding the physics modules changes \iou{} by less than
seed noise: all variants---including a no-physics backbone-plus-decoder---lie in
$[0.9317,0.9350]$ ($\Delta=0.0033$), while the seed standard deviation is
$0.0008$, and a paired bootstrap confirms the largest physics effect
($+0.003$ \iou) is practically equivalent to zero within a $\pm0.01$ margin
(Sec.~\ref{sec:ablation}). A real-video control (Sec.~\ref{sec:cross})
shows the learned \ofcv{} gate is inert even given genuine parallax. The
accuracy is a property of the representation: a \dinov{} backbone with a plain
decoder (the 30-epoch ``neither'' variant) reaches $0.932$ \iou, against the
10-epoch \spectra{} configuration at $0.9217$.

\textbf{On the 0.92 number, stated plainly.} \spectra{} reaches a competitive
binary foreground \iou{} of $0.9217$ (val) / $0.9237$ (test), above
ImageNet-pretrained CNN baselines under a matched 10-epoch budget
(DeepLabV3+ $0.8856$). We do \emph{not} claim this as a win for our method. The
comparison contrasts a \dinov{} self-supervised ViT (pretrained on the large
LVD-142M corpus) against ImageNet-1k--pretrained CNN encoders fine-tuned for the
same number of epochs; it therefore conflates \emph{backbone and pretraining
corpus} with anything we contribute, and the ablation shows the contribution is
null. We present this number purely to establish that the operating point is
reasonable (and that zero-shot SAM, at $0.1341$, is not), and we caveat it as
such everywhere it appears (Sec.~\ref{sec:results}).

\textbf{Contributions.}
\begin{enumerate}
\item \textbf{A controlled refutation.} We isolate two representative
physics-motivated cues and show, across five variants, 30 epochs, and three
seeds, that they add no accuracy beyond seed noise on Trans10K
(Sec.~\ref{sec:ablation}).
\item \textbf{A real-video falsification.} On the VGSD video benchmark, supplying
genuine inter-frame parallax changes the \ofcv{} map by $<0.05\%$ and the
prediction by $<0.001\%$ versus a duplicated frame---the flow cue is empirically
blind to motion (Sec.~\ref{sec:cross}).
\item \textbf{A representation finding.} The measured accuracy is attributable to
the self-supervised backbone, not the cues or the architecture; we make the
backbone/pretraining confound in the baseline comparison explicit rather than
exploiting it (Secs.~\ref{sec:results},~\ref{sec:ablation}).
\item \textbf{Supporting analyses.} Zero-shot cross-dataset transfer (ClearPose,
VGSD), robustness over eight corruptions $\times$ five severities, an automated
failure taxonomy over all $4{,}428$ test images, and an interpretability
analysis dissociating interpretable cues from accuracy
(Secs.~\ref{sec:cross}--\ref{sec:interp}).
\end{enumerate}

\textbf{Why this matters.} Physics-motivated and ``causal'' inductive biases are
increasingly proposed for ill-posed perception tasks, often validated only by a
headline metric on one dataset against weaker backbones. We show, on a task
where such a prior is especially intuitive, that a popular cue can be
\emph{interpretable, non-collapsing, and entirely inert}---and that an
unexamined backbone advantage can masquerade as a method effect. Reporting this
with full controls helps the community avoid a recurring attribution error and
sets a reproducible template for testing, rather than assuming, that a physical
prior contributes.

We deliberately do \emph{not} compare against published Trans10K leaderboard
numbers: they are reported under the multi-class mean-\iou{} protocol, whereas we
report binary foreground \iou, so the two are not comparable
(Sec.~\ref{sec:results}). All numbers herein are our own, measured under one
fixed pipeline.

%==============================================================================
